\subsection{The classical class-number interpretation of the local criterion}
\label{subsec:cm-class-number}
The determinant in the local criterion is governed by the classical
relative class-number formula.  We record the normalization explicitly
to distinguish invertibility of this integral operator from the later
question of global arithmetic saturation.

Let \(p\geq5\) be prime, put \(h=(p-1)/2\), choose a primitive residue
\(g\) modulo \(p\), and let \(s=[g]\), \(c=s^h\).  As above,
\[
 \Phi^\ast=\sum_{a=1}^{h}[a]^{-1},\qquad
 (s-1)\Phi^\ast=(c-1)Q_p,\qquad
 R_p=\mathbf Z[s]/(s^h+1).
\]
Let \(h_p^-\) be the relative class number of
\(\mathbf Q(\zeta_p)\), let \(d=\operatorname{ord}_p(2)\), and put
\(r=(p-1)/d\).

\begin{proposition}\label{prop:cm-class-number-determinant}
Multiplication by \(Q_p\) on \(R_p\) has positive determinant
\[
 \det(Q_p\mid R_p)=\frac{h_p^-}{p}
 \begin{cases}
 (2^{d/2}+1)^r,&d\ \text{even},\\
 (2^d-1)^{r/2},&d\ \text{odd}.
 \end{cases}
\]
Consequently the determinant is nonzero for every \(p\), and
\[
 v_2\bigl(\det(Q_p\mid R_p)\bigr)=v_2(h_p^-).
\]
The operator \(Q_p\) is a unit in
\(\mathbf Z_2[s]/(s^h+1)\) if and only if \(h_p^-\) is odd,
equivalently if and only if the class number of
\(\mathbf Q(\zeta_p)\) is odd.
\end{proposition}
\begin{proof}
For an odd character \(\psi\) of
\((\mathbf Z/p)^\times\), set
\(B_{1,\psi}=p^{-1}\sum_{a=1}^{p-1}a\psi(a)\).
Permutation of the nonzero residues by multiplication by \(2\),
using the least positive representatives and
\(\sum_{a>h}\psi(a)=-\sum_{a\leq h}\psi(a)\), gives
\[
 \sum_{a=1}^{h}\psi(a)
       =(\psi(2)^{-1}-2)B_{1,\psi}.
\]
Evaluation of the defining identity at an odd character \(\chi\)
therefore gives
\[
 Q_p(\chi)=
 \frac{(1-\chi(g))(\chi(2)-2)B_{1,\chi^{-1}}}{2}.
\]
The odd values of \(\chi(g)\) are the roots of \(X^h+1\), so
their factors \(1-\chi(g)\) have product \(2\).  Kummer's
relative class-number formula, in the normalization
\cite[Eq.~(1.2), p.~1]{Agoh2012}, is
\[
 h_p^-=2p\prod_{\chi\ {\rm odd}}\left(-\frac{B_{1,\chi}}2\right).
\]
Multiplying the evaluated identities yields
\[
 \det(Q_p\mid R_p)
     =\frac{h_p^-}{p}\prod_{\chi\ {\rm odd}}(2-\chi(2)).
\]

If \(d\) is even, then \(2^{d/2}=-1\pmod p\), and the values
\(\chi(2)\) for odd \(\chi\) are all roots of \(X^{d/2}+1\),
each with multiplicity \(r\).  If \(d\) is odd, they are all
roots of \(X^d-1\), each with multiplicity \(r/2\).  This gives
the displayed positive formula.  Its factors other than
\(h_p^-\) are odd in \(\mathbf Z_2^\times\), proving the valuation
identity.  An element of a finite free \(\mathbf Z_2\)-algebra
is a unit precisely when its multiplication determinant is
a unit: the inverse linear map supplies a multiplicative inverse
by its value on \(1\).  Finally the equivalence of odd relative
and odd full cyclotomic class number is
\cite[Proposition~1, equivalence (1)--(2), p.~207]{Dummit2020}.
\end{proof}

Thus the prime set \(\Pi\) in the simultaneous local
normalization theorem can equivalently be written
\[
 \Pi=\{p\geq5:\ h(\mathbf Q(\zeta_p))\text{ is odd}\}.
\]
This criterion concerns the local reflex-norm operator.
The degree-\(29\) and degree-\(113\) global arguments show that
an even class number need not prevent arithmetic saturation.
When the local index is computed on the unit group itself,
the finite sign subgroup in the integral index formula must
still be included.
